\centerline{\bf \Huge Памятка по неравенствам}

\begin{flushright}
        \scriptsize{я устал, босс...}
\end{flushright}

\textbf{\Large Неравенства о средних.}

\textbf{I. Теория}

Выражения
$$
\sqrt{\frac{x_1^2+x_2^2+\ldots+x_n^2}{n}} \geqslant \frac{x_1+x_2+\ldots+x_n}{n} \geqslant \sqrt[n]{x_1 x_2 \ldots x_n} \geqslant \frac{n}{\frac{1}{x_1}+\frac{1}{x_2}+\ldots+\frac{1}{x_n}}
$$


называются соответственно средним квадратическим, средним арифметическим, средним геометрическим и средним гармоническим неотрицательных чисел $x_1, x_2, \ldots, x_n$ и удовлетворяют предлагаемому неравенству.

Частный случай этого неравенства при $n=2$, который часто называют \textit{неравенством Коши}, в особенности бывает очень полезен. 

\textbf{II Примеры использования}

\textit{Пример 1.}
$\bigl[\textbf{Регион ВсОШ 1996, 10.1}\bigr]$
Докажите, что если $a, b, c$ - положительные числа и $a b+b c+c a>a+b+c$, то $a+b+c>3$.

\textbf{Решение.}
Используем формулу сокращенного умножения:

$$
(a+b+c)^2=a^2+b^2+c^2+2 a b+2 b c+2 a c
$$


По неравенству Коши имеем:

$$
a^2+b^2+c^2=\frac{a^2+b^2}{2}+\frac{b^2+c^2}{2}+\frac{c^2+a^2}{2} \geqslant a b+b c+c a
$$


Таким образом, $(a+b+c)^2 \geqslant 3(a b+b c+a c)>3(a+b+c)$, откуда и получаем $a+b+c>3$.

\textit{Пример 2.}
$\bigl[\textbf{Регион ВсОШ 2025, 11.2}\bigr]$
Вещественные числа $x, y, z$ таковы, что $2 x>y^2+z^2, 2 y>z^2+x^2, 2 z>x^2+y^2$. Докажите, что каждое из чисел $x, y, z$ меньше 1 .

\textbf{Решение.}
Третье решение. Из условия следует, что $2 x>y^2+z^2 \geqslant 0$, аналогично $y, z>0$. Также $2 x>y^2+z^2 \geqslant 2 y z$ по неравенству о средних. Значит, $x>y z$, аналогично $y>z x$ и $z>x y$. Не умаляя общности можно считать, что $x$ - минимальное из чисел $x, y, z$, тогда $y \geqslant x>y z$, откуда $z<1$, аналогично $y<1$, а тогда и $x<1$.

\textbf{III Гринфлаги}

В целом неравенства о средних можно использовать хоть когда, это универсальный способ решения неравенств. В первую очередь обращайте внимание, нет ли в задаче суммы чисел ($x+y+z...$), суммы квадратов ($x^2+y^2+z^2...$) и тому подобное. То есть всякое выражение, которое напоминало бы нам неравенство о средних. Но даже если не напоминает, можно пробовать <<усреднять>> любое выражение, авось прокатит.

\textbf{\Large Неравенство КБШ и неравенство Седракяна.}

\textbf{I Теория}

\textit{Неравенство Коши-Буняковского-Шварца.}
Пусть $a_1, a_2, \ldots, a_n$ и $b_1, b_2, \ldots, b_n$ - это действительные числа. Тогда справедливо неравенство:

$$
\left(a_1^2+a_2^2+\cdots+a_n^2\right)\left(b_1^2+b_2^2+\cdots+b_n^2\right) \geqslant\left(a_1 b_1+a_2 b_2+\cdots+a_n b_n\right)^2 .
$$

\textbf{II Примеры использования}

\textit{Пример 1.} 
Зачастую встречаются в неравенствах, где слева стоит константа, умноженная на какое-то выражение, а справа просто выражение. Например, нам нужно доказать вот такое неравенство:

$$
    n(a_1^2 + a_2^2 + \ldots a_n^2) \geqslant (a_1 + a_2 + \ldots + a_n)^2
$$

Вот эту константу $n$ удобно бывает представить как $n = 1^2 + 1^2 + \ldots + 1$ --- сумма $n$ единичек. Тогда по неравенству КБШ получаем:

$$
    \left(a_1^2+a_2^2+\cdots+a_n^2\right)\left(1^2+1^2+\cdots+1^2\right) \geqslant\left(a_1 \cdot 1+a_2 \cdot 1+\cdots+a_n \cdot 1\right)^2
$$

\textit{Пример 2.}
Если в условии уже дано что-то про сумму квадратов, это хороший повод попробовать КБШ.

Пусть $a^2+b^2+c^2=2$. Докажите неравенство $a+b+c \leqslant a b c+2$.

Поймём для начала что-то про $b c$. Перепишем равенство из условия в виде $b^2+c^2=2-a^2$. К тому же верна следующая цепочка неравенств

$$
2 b c \leqslant b^2+c^2=2-a^2 \leqslant 2
$$


Значит, получаем, что $b c \leqslant 1$. Перенесём теперь $a b c$ в левую сторону и запишем КБШ

$$
b+c+a(1-b c) \leqslant \sqrt{\left((b+c)^2+a^2\right)\left(1^2+(1-b c)^2\right)}=\sqrt{(2+2 b c)\left(1+(1-b c)^2\right)}
$$


Получаем, что нам надо доказать следующее неравенство $\sqrt{(2+2 b c)\left(1+(1-b c)^2\right)} \leqslant 2$. Возведём в квадрат, сделаем замену $1-b c=x$, где $x$ неотрицательный, и сделаем преобразования

$$
\begin{gathered}
(4-2 x)\left(1+x^2\right) \leqslant 4 \\
4 x^2-2 x-2 x^3 \leqslant 0 \\
-2 x(x-1)^2 \leqslant 0
\end{gathered}
$$


Последнее неравенство верно, поэтому получаем, что и наше исходное неравенство доказано.

\textit{Пример 3. (Геометрический)}

Если говорить русским языком, то \textbf{произведение длин векторов всегда не меньше чем их скалярное произведение причём в пространстве любой размерности.}

Даны числа $x, y, z$ такие, что

$$
4^x+\sin ^4 y+\ln ^6 z=16
$$


Докажите, что

$$
2^{x+1}+3 \sin ^2 y-6 \ln ^3 z \leqslant 28
$$


Используем неравенство КБШ в векторном виде.

Рассмотрим векторы $\vec{a}=\left(2^x ; \sin ^2 y ; \ln ^3 z\right)$ и $\vec{b}=(2 ; 3 ;-6)$. Скалярное произведение

$$
\vec{a} \cdot \vec{b}=2^{x+1}+3 \sin ^2 y-6 \ln ^3 z \leqslant|\vec{a}| \cdot|\vec{b}|
$$


Имеем

$$
\begin{gathered}
|\vec{b}|=\sqrt{4+9+36}=7 \\
|\vec{a}|=\sqrt{4^x+\sin ^4 y+\ln ^6 z}=4
\end{gathered}
$$


Тогда получаем, что

$$
2^{x+1}+3 \sin ^2 y-6 \ln ^3 z \leqslant 28
$$

\textit{Неравенство Седракяна (дробный КБШ).}
Пусть $a_1, a_2, \ldots, a_n$ и $b_1, b_2, \ldots, b_n$ - это действительные положительные числа, тогда справедливо неравенство:

$$
\dfrac{a_1^2}{b_1}+\dfrac{a_2^2}{b_2}+\cdots+\dfrac{a_n^2}{b_n} \geqslant \dfrac{\left(a_1+a_2+\cdots+a_n\right)^2}{b_1+b_2+\cdots+b_n}
$$

Равенство приходится тогда и только тогда, когда $\dfrac{a_1}{b_1}=\dfrac{a_2}{b_2}=\cdots=\dfrac{a_n}{b_n}$.

\textit{Пример 1.}
Числа $a, b, c, d$ положительны и причём $a+b+c+d=1$. Найдите наименьшее возможное значение выражения

$$
\dfrac{a^2}{1-a}+\dfrac{b^2}{1-b}+\dfrac{c^2}{1-c}+\dfrac{d^2}{1-d}
$$

По неравенству Седракяна сразу же получаем

$$
\dfrac{a^2}{1-a}+\dfrac{b^2}{1-b}+\dfrac{c^2}{1-c}+\dfrac{d^2}{1-d} \geqslant \dfrac{(a+b+c+d)^2}{1-a+1-b+1-c+1-d}=\dfrac{1}{3}
$$

Причём равенство достигается, когда $a = b = c = d = \dfrac{1}{4}$

\textit{Пример 2.}
Про положительные числа $a, b, c$ известно, что $\dfrac{1}{a}+\dfrac{1}{b}+\dfrac{1}{c}=6$. Докажите неравенство

$$
\dfrac{a^2+b}{a+b}+\dfrac{b^2+c}{b+c}+\dfrac{c^2+a}{c+a} \geqslant \dfrac{9}{4}
$$

Запишем левую часть в следующем виде, и оценим сначала первую сумму, а потом сумму, которую вычиатем.

$$
\dfrac{a^2+\dfrac{1}{4}+b}{a+b}+\dfrac{b^2+\dfrac{1}{4}+c}{b+c}+\dfrac{c^2+\dfrac{1}{4}+a}{c+a}-\left(\dfrac{1}{4(a+b)}+\dfrac{1}{4(b+c)}+\dfrac{1}{4(a+c)}\right)
$$


Числитель этих дробей можно оценить с помощью следующего неравенства $a^2+\dfrac{1}{4} \geqslant a$, откуда получим такую оценку суммы:

$$
\dfrac{a^2+\dfrac{1}{4}+b}{a+b}+\dfrac{b^2+\dfrac{1}{4}+c}{b+c}+\dfrac{c^2+\dfrac{1}{4}+a}{c+a} \geqslant \dfrac{a+b}{a+b}+\dfrac{b+c}{b+c}+\dfrac{c+a}{c+a}=3
$$


А вычитаемую сумму запишем таким образом:

$$
\dfrac{\left(\dfrac{1}{4}+\dfrac{1}{4}\right)^2}{(a+b)}+\dfrac{\left(\dfrac{1}{4}+\dfrac{1}{4}\right)^2}{(b+c)}+\dfrac{\left(\dfrac{1}{4}+\dfrac{1}{4}\right)^2}{(a+c)}
$$


Мы можем оценить каждую из дробей по неравенству КБШ для дробей, откуда получим следующее

$$
\dfrac{\left(\dfrac{1}{4}+\dfrac{1}{4}\right)^2}{(a+b)}+\dfrac{\left(\dfrac{1}{4}+\dfrac{1}{4}\right)^2}{(b+c)}+\dfrac{\left(\dfrac{1}{4}+\dfrac{1}{4}\right)^2}{(a+c)} \leqslant\left(\dfrac{\dfrac{1}{4^2}}{a}+\dfrac{\dfrac{1}{4^2}}{b}\right)+\left(\dfrac{\dfrac{1}{4^2}}{b}+\dfrac{\dfrac{1}{4^2}}{c}\right)+\left(\dfrac{\dfrac{1}{4^2}}{c}+\dfrac{\dfrac{1}{4^2}}{a}\right)=
$$

$$
=\dfrac{1}{8} \cdot\left(\dfrac{1}{a}+\dfrac{1}{b}+\dfrac{1}{c}\right)=\dfrac{3}{4}
$$


В итоге получаем, что и требовалось доказать

$$
\begin{gathered}
\dfrac{a^2+b}{a+b}+\dfrac{b^2+c}{b+c}+\dfrac{c^2+a}{c+a}= \\
=\dfrac{a^2+\dfrac{1}{4}+b}{a+b}+\dfrac{b^2+\dfrac{1}{4}+c}{b+c}+\dfrac{c^2+\dfrac{1}{4}+a}{c+a}-\left(\dfrac{1}{4(a+b)}+\dfrac{1}{4(b+c)}+\dfrac{1}{4(a+c)}\right) \geqslant 3-\dfrac{3}{4}=\dfrac{9}{4}
\end{gathered}
$$

\textit{Пример 3.}
Для положительных $a, b, c, d$ докажите неравенство

$$
\frac{12}{a+b+c+d} \leqslant \frac{1}{a+b}+\frac{1}{a+c}+\frac{1}{a+d}+\frac{1}{b+c}+\frac{1}{b+d}+\frac{1}{c+d} \leqslant \frac{3}{4}\left(\frac{1}{a}+\frac{1}{b}+\frac{1}{c}+\frac{1}{d}\right)
$$

Воспользуемся КБШ для дробей:

$$
\begin{aligned}
\frac{1}{a+b}+\frac{1}{a+c}+\frac{1}{a+d}+\frac{1}{b+c} & +\frac{1}{b+d}+\frac{1}{c+d} \geqslant \frac{(1+1+1+1+1+1)^2}{3(a+b+c+d)}= \\
& =\frac{12}{a+b+c+d}
\end{aligned}
$$


Воспользуемся КБШ для дробей для правой части следующим образом:

$$
\begin{gathered}
\frac{3}{4}\left(\frac{1}{a}+\frac{1}{b}+\frac{1}{c}+\frac{1}{d}\right)=\frac{1}{4}\left(\frac{3}{a}+\frac{3}{b}+\frac{3}{c}+\frac{3}{d}\right)= \\
=\frac{3}{4}\left(\left(\frac{1}{a}+\frac{1}{b}\right)+\left(\frac{1}{a}+\frac{1}{c}\right)+\left(\frac{1}{a}+\frac{1}{d}\right)+\left(\frac{1}{b}+\frac{1}{c}\right)+\left(\frac{1}{b}+\frac{1}{d}\right)+\left(\frac{1}{c}+\frac{1}{d}\right)\right) \geqslant \\
\geqslant \frac{1}{4}\left(\frac{4}{a+b}+\frac{4}{a+c}+\frac{4}{a+d}+\frac{4}{b+c}+\frac{4}{b+d}+\frac{4}{c+d}\right)= \\
=\frac{1}{a+b}+\frac{1}{a+c}+\frac{1}{a+d}+\frac{1}{b+c}+\frac{1}{b+d}+\frac{1}{c+d}
\end{gathered}
$$

\textbf{III Гринфлаги}

На самом то деле это понятно из условия самих неравенств. КБШ обычный применяем, когда видим хорошую сумму чисел в каких-то степенях, будь то просто сумма или сумма квадратов, кубов и т.д. В формулировке фигурируют два набора чисел и их соответствующие произведение. Если вы видите два набора и у них хорошие произведение чисел с одинаковыми номерами --- смело делайте КБШ. Так же, если у вас в условии дано, что сумма фиксирована помимо тригера к методу Штурма должен срабатывать и триггер к неравенству КБШ. В случае с дробным всё проще. Всё то же самое, что и с обычным КБШ, но ещё к этому добавляется, что его \textbf{ВСЕГДА} можно использовать для суммы \textbf{положительных} дробей. Даже если у вас в числителях нет квадратов, вы сами можете их организовать просто взяв набор из корней этих числителей.

Ещё к примеру, можно домножить числитель и знаменатель на что-нибудь или ревёрснуть: вместо дробей $\frac{P}{Q}$ рассмотреть дроби $R-\frac{P}{Q}=\frac{R Q-P}{Q}$, где $P, Q, R$ - некоторые выражения (смысл в том, что новые дроби надо оценить с другой стороны).

\textbf{\Large Неравенство Мюрхеда и Шура.}

\textbf{I Теория}

Пусть даны переменные $x_1, x_2, \ldots, x_n$. Однородным симметрическим многочленом, порождённым упорядоченным набором целых неотрицательных чисел $a_1 \geqslant a_2 \geqslant \ldots \geqslant a_n$, называют выражение $T_{a_1, a_2, \ldots, a_n}=\sum_{\sigma \in S_n} x_{\sigma(1)}^{a_1} x_{\sigma(2)}^{a_2} \ldots x_{\sigma(n)}^{a_n}$, где суммирование ведётся по множеству $S_n$ всех перестановок множества $\{1,2, \ldots, n\}$.

Пусть даны два упорядоченных набора целых неотрицательных чисел $a_1 \geqslant a_2 \geqslant \ldots \geqslant a_n$ и $b_1 \geqslant b_2 \geqslant \geqslant \ldots \geqslant b_n$. Тогда набор $a_k$ \textit{мажорирует} набор $b_k$, если выполнены следующие условия: $a_1 \geqslant b_1, a_1+ +a_2 \geqslant b_1+b_2, \ldots, a_1+\ldots+a_{n-1} \geqslant b_1+\ldots+b_{n-1}, a_1+\ldots+a_n=b_1+\ldots+b_n$. Если набор $a_k$ мажорирует набор $b_k$, то для неотрицательных значений переменных верно \textit{неравенство Мюрхеда}: $T_{a_1, a_2, \ldots, a_n} \geqslant T_{b_1, b_2, \ldots, b_n}$.


\textit{Неравенство Шура.}
Для любого натурального $n$ и любых неотрицательных значений переменных верно, что $T_{n+2,0,0}+T_{n, 1,1} \geqslant 2 T_{n+1,1,0}$. 


\textbf{II Примеры использования}

\textit{Пример 1.}
Для положительных $a, b, c$ докажите неравенство

$$
\frac{a}{b+c}+\frac{b}{c+a}+\frac{c}{a+b} \geqslant \frac{3}{2}
$$

\textbf{Решение.}
Умножим неравенство на все встречающиеся знаменатели и раскроем все скобки. Получившееся представляем в виде выражений $T_{i, j, k}(a, b, c)$ из неравенства Мюрхеда

$$
\begin{gathered}
T_{3,0,0}(a, b, c)+2 T_{2,1,0}(a, b, c)+T_{1,1,1}(a, b, c) \geqslant T_{1,1,1}(a, b, c)+3 T_{2,1,0}(a, b, c) \\
T_{3,0,0}(a, b, c) \geqslant T_{2,1,0}(a, b, c)
\end{gathered}
$$


Последнее неравенство верно по неравенству Мюрхеда.

\textit{Пример 2.}
Для положительных чисел $a, b, c$ докажите неравенство

$$
\frac{(a+b-c)^2}{a b}+\frac{(b+c-a)^2}{b c}+\frac{(c+a-b)^2}{c a} \geqslant 3
$$

Домножим на $a b c$, раскроем скобки, после этого имеем:

$$
T_{3,0,0}(a, b, c)+T_{1,1,1}(a, b, c) \geqslant 2 T_{2,1,0}(a, b, c)
$$


что очевидно по неравенству Шура.

\textit{Пример 3.}
Для положительных $a, b, c, d$ докажите неравенство:

$$
\sqrt{\frac{a^2+b^2+c^2+d^2}{4}} \geqslant \sqrt[3]{\frac{a b c+b c d+c d a+d a b}{4}} .
$$


\textbf{Решение.}
Возведём в шестую степень. Теперь достаточно доказать, что

$$
\left(a^2+b^2+c^2+d^2\right)^3 \geqslant 4 \cdot(a b c+b c d+c d a+d a b)^2 .
$$


После раскрытия скобок останется следующее неравенство:

$$
\frac{1}{6} \cdot T(6,0,0,0)+\frac{3}{2} \cdot T(4,2,0,0)+T(2,2,2,0) \geqslant 4 \cdot\left(\frac{1}{6} \cdot T(2,2,2,0)+\frac{1}{2} \cdot T(2,2,1,1)\right) .
$$


Оно верно, так как суммы коэффициентов перед $T$ одинаковые, а все наборы слева мажорируют наборы справа.

\textit{Комментарий: можно было не раскрывать скобки, а просто прикинуть, какие наборы параметров получаются слева и справа. И так понятно, что минимальный по мажорированию набор слева - это (2, 2, 2, 0), а максимальный справа - тоже (2, 2, 2, 0). Суммы коэффициентов должны получиться одинаковыми, так как при $a=b=c=d$ наблюдается равенство, а значит неравенство, возникающее при раскрытии скобок, может быть получено как линейная комбинация нескольких неравенств Мюрхеда.}

\textbf{III Гринфлаги}

Видите симметрическое выражение? Поздравляю, вы решите эту задачу с помощью неравенства Мюрхеда/Шура. Нет симметрического выражения? Разбейте его на сумму симметрических и оцените отдельно каждое по Мюрхеду/Шура (пример 3). Не получается ни то ни другое? Ну печально, что могу сказать.

\textbf{\Large Гомогенизация.}

\textbf{I Теория}

Введение дополнительного параметра $\lambda$ и замена переменных вида $a=\frac{a^{\prime}}{\lambda}$ позволяет свести симметрическое неоднородное неравенство со связкой к симметрическому однородному без связки. (Достаточно параметр $\lambda$ выразить в связке и в неравенстве и сопоставить полученные оценки). Обычно после замены штрихи не пишут для удобства. 

\textbf{II Примеры использования}

\textit{Пример 1}
Положительные числа $a, b, c$ таковы, что $a^2+b^2+c^2=3$. Докажите:

$$
a+b+c \geqslant a b+b c+c a .
$$


\textbf{Решение.}
Сделаем замену $a=a^{\prime} / \lambda, b=b^{\prime} / \lambda, c=c^{\prime} / \lambda$. Получим следующую задачу (далее штрихи не пишем).
Положительные числа $a, b, c, \lambda$ таковы, что $a^2+b^2+c^2=3 \lambda^2$. Докажите: $\lambda \cdot(a+b+c) \geqslant a b+b c+c a$.
Выразим $\lambda$ из связки и подставим в желаемое неравенство. Докажем, что при всех положительных $a, b, c$ верно

$$
\sqrt{\frac{a^2+b^2+c^2}{3}} \cdot(a+b+c) \geqslant(a b+b c+c a) .
$$


После возведения в квадрат получим $\left(a^2+b^2+c^2\right) \cdot(a+b+c)^2 \geqslant 3 \cdot(a b+b c+c a)^2$. Раскроем скобки:

$$
\frac{1}{2} \cdot T(4,0,0)+T(2,2,0)+2 \cdot T(3,1,0)+T(2,1,1) \geqslant \frac{3}{2} \cdot T(2,2,0)+3 \cdot T(2,1,1) .
$$


Чтобы получить последнее, достаточно сложить несколько неравенств Мюрхеда и равенств, а именно:

$$
\frac{1}{2} \cdot T(4,0,0) \geqslant \frac{1}{2} \cdot T(2,2,0) ; T(2,2,0)=T(2,2,0) ; 2 \cdot T(3,1,0) \geqslant 2 \cdot T(2,1,1) ; T(2,1,1)=T(2,1,1) .
$$

\textit{Пример 2.}
Для положительных $a, b, c$ известно, что $\frac{3}{a b c} \geqslant a+b+c$. Покажите, что выполнено неравенство

$$
\frac{1}{a}+\frac{1}{b}+\frac{1}{c} \geqslant a+b+c
$$

\textbf{Решение.} После гомогенизации (замены из решения предыдущей задачи) получим вот что: Для положительных $a, b, c, \lambda$ известно, что $3 \cdot \lambda^4 \geqslant(a+b+c) \cdot a b c$. Покажите, что $\lambda^2 \cdot(a b+b c+c a) \geqslant(a+b+c) \cdot a b c$. В общем, достаточно для положительных $a, b, c$ доказать неравенство:

$$
(a b+b c+c a) \cdot \sqrt{\frac{(a+b+c) \cdot a b c}{3}} \geqslant(a+b+c) \cdot a b c .
$$

Возведём в квадрат, сократим, перепишем: $(a b+b c+c a)^2 \geqslant 3 a b c \cdot(a+b+c)$. После раскрытия скобок очевидно.

\textbf{III Гринфлаги}

Суть гомогенизации в том, что после неё ВСЕГДА будет получено однородное/симметричное неравенство. Полученное однородное неравенство или решается втупую, или сводится к неравенству Мюрхеда, или сводится к неравенствами Мюрхеда и Шура, или получается слишком громоздким и сложным, и тогда от идеи гомогенизации в качестве первого шага решения стоит отказаться. Использовать гомогенизацию стоит в том случае, когда у вас дано в условии одно неравенство, а нужно доказать другое.

\textbf{\Large Выпуклость.}

\textbf{I. Теория}

\textbf{Определение.} Функция $f: \mathbb{R} \rightarrow \mathbb{R}$ называется \textit{выпуклой} (или же \textit{выпуклой вниз}), если для любых точек $X, Y$ из $\mathbb{R}$ и любых вещественных неотрицательных $\lambda$ и $\mu$, удовлетворяющих $\lambda+\mu=$ 1 , верно неравенство: $\lambda f(X)+\mu f(Y) \geqslant f(\lambda X+\mu Y)$. Геометрический смысл написанного: хорда, стягивающая любые две точки графика функции, лежит над графиком.

\textbf{Полезный факт 1.} Выпуклая вниз (вверх) на отрезке $[a,b]$ функция принимает максимальное (минимальное) значение в одном из краёв отрезка.

\textbf{Полезный факт 2.} Сумма/разность выпуклых вверх (вниз) функций будет выпуклой вверх (вниз) функцией.

\textbf{II. Пример использования}

Для $a, b, c \in[0,1]$ докажите неравенство

$$
\frac{a}{b+c+1}+\frac{b}{c+a+1}+\frac{c}{a+b+1}+(1-a)(1-b)(1-c) \leqslant 1 .
$$

\textbf{Решение.}
Рассмотрим выражение слева, как функцию от переменной $a \in [0,1]$. Заметим, что каждое слагаемое слева является выпуклой вниз функцией, а значит всё выражение является выпуклой вниз функцией. Тогда максимальное значение принимается либо в точке $a=0$, либо в точке $a=1$. Аналогично поступаем с остальными переменными и нам лишь остается проверить, что выражение слева при наборах $(0,0,0), (0,0,1), (0,1,1), (1,1,1)$ будет не больше 1. 

\textbf{III Грин флаги}

Тут надо чувствовать.